\documentclass{article}
\usepackage{amsmath, amssymb, amsthm}

\setlength{\parindent}{0pt}

\newtheorem{claim}{Claim}

\begin{document}

\begin{claim}
    If a Harvard BS graduate starts with an annual salary of $\$140{,}000$ and receives
    a guaranteed $\$25{,}000$ raise each year, while an MIT SB graduate starts with
    $\$100{,}000$ and receives a guaranteed $15\%$ raise each year, then after working
    the same number of years, the MIT SB graduate will eventually earn more in total.
\end{claim}

\begin{proof}
    Let $n$ denote the number of years both graduates work. The Harvard BS graduate's
    annual salary grows linearly, so the total earnings after $n$ years are $O(n^2)$.

    \medskip

    The MIT SB graduate's annual salary grows exponentially. Hence the total earnings
    after $n$ years are at least as large as the final year's salary, i.e. 
    \[
    100000(1.15)^{n - 1} = \Theta((1.15)^n)
    \]

    \medskip

    Since
    \[
        n^2 = o((1.15)^n)
    \]
    exponential growth eventually dominates quadratic growth. Thus, for
    sufficiently large $n$, the MIT SB graduate's total earnings exceed those of the
    Harvard BS graduate.
\end{proof}

\begin{claim}
    Assuming the bankrate is fixed at $3\%$ per year, the MIT graduate's salary package 
    is better than the Harvard graduate's salary package when $n \geq 16$. 
\end{claim}

\begin{proof}
    Let
    \[
    H_n = \sum_{k=0}^{n-1}\frac{140000 + 25000k}{1.03^k}
    \]
    and
    \[
    M_n = \sum_{k=0}^{n-1}\frac{100000(1.15)^k}{1.03^k}
    \]
    denote the present values of the Harvard and MIT salary packages respectively and 
    define $D_n = H_n - M_n$.

    \medskip

    Using the closed forms for a finite geometric series and a finite arithmetic-geometric 
    series, we compute
    \[
    D_{15} > 0 \quad \text{and} \quad D_{16} < 0
    \]

    \medskip

    Now we show that $D_n$ is strictly decreasing for all $n \ge 16$. Since
    \[
    D_{n+1} - D_n = \frac{140000 + 25000n - 100000(1.15)^n}{1.03^n}
    \]
    it suffices to prove
    \[
    100000(1.15)^n > 140000 + 25000n
    \]
    for every $n \ge 16$.

    \medskip

    This inequality holds when $n = 16$. Moreover, if it holds for some $n\ge16$, then
    \[
    \begin{aligned}
    100000(1.15)^{n + 1} &= 1.15 \cdot 100000(1.15)^n \\
    &> 1.15(140000 + 25000n) \\
    &= 161000 + 28750n \\
    &> 165000 + 25000n \\
    &= 140000 + 25000(n + 1)
    \end{aligned}
    \]
    since
    \[
    (161000 + 28750n) - (165000 + 25000n) = 3750n - 4000 > 0
    \]
    for every $n \ge 16$. Hence, by induction,
    \[
    100000(1.15)^n > 140000 + 25000n
    \]
    for all $n \ge 16$. Therefore $D_n$ is strictly decreasing for all $n \ge 16$. 
    Since $D_{16} < 0$, it follows that $D_n < 0$ for all $n \ge 16$.

    \medskip

    Hence the MIT graduate's salary package is better than the Harvard graduate's salary 
    package when $n \geq 16$.
\end{proof}

\end{document}
